\begin{description}
\item[(T3.1)]
\[
\frac{\rho_{\mathrm{m_0}}/\rho_{\mathrm{c}}}{\rho_{\mathrm{r_0}}/\rho_{\mathrm{c}}}
= \frac{\Omega_{\mathrm{m_0}}}{\Omega_{\mathrm{r_0}}}
= \frac{0.3}{10^{-4}} = 3000
\]
At $z_{\mathrm{e}}$, both matter density and radiation density were equal.
\begin{align*}
\rho_{\mathrm{r}} &= \rho_{\mathrm{m}}\\
\therefore\ \rho_{\mathrm{r_0}}(1+z_{\mathrm{e}})^4
  &= \rho_{\mathrm{m_0}}(1+z_{\mathrm{e}})^3\\
1+z_{\mathrm{e}} &= \frac{\rho_{\mathrm{m_0}}}{\rho_{\mathrm{r_0}}} = 3000
  \tag*{(2.0)}\\
\therefore\ z_{\mathrm{e}} &\simeq 3000 \tag*{(1.0)}
\end{align*}
\textbf{Only $z_{\mathrm{e}} = 2999$ and $z_{\mathrm{e}} = 3000$ are
acceptable answers.}

\item[(T3.2)] As the Universe behaves like an ideal black body, the radiation
density will be proportional to the fourth power of the temperature (Stefan's
law).
\begin{align*}
\left(\frac{T_{\mathrm{e}}}{T_0}\right)^4
  &= \frac{\rho_{\mathrm{r_e}}}{\rho_{\mathrm{r_0}}}
   = \frac{\rho_{\mathrm{r_0}}(1+z_{\mathrm{e}})^4}{\rho_{\mathrm{r_0}}}
  \tag*{(2.0)}\\
\left(\frac{T_{\mathrm{e}}}{2.732}\right)^4
  &= (1+z_{\mathrm{e}})^4 \tag*{(1.0)}\\
\frac{T_{\mathrm{e}}}{2.732} &= 1+z_{\mathrm{e}} = 3000\\
T_{\mathrm{e}} &= 3000\times2.732\\
T_{\mathrm{e}} &= 8200\ \mathrm{K} \tag*{(1.0)}
\end{align*}
\textbf{$8100 \le T_{\mathrm{e}} \le 8200$ gives 1.0;
$8200 < T_{\mathrm{e}} \le 9000$ gives 0.5; else 0.}

\item[(T3.3)] Wien's law:
\begin{align*}
\lambda_{\max} &= \frac{0.002898\ \mathrm{m\,K}}{T_e}\\
  &= \frac{0.002898}{8200}\ \mathrm{m} = 354\ \mathrm{nm} \tag*{(1.0)}\\
E_\nu &= \frac{hc}{\lambda_{\max}}\\
  &= \frac{6.62\times10^{-34}\times3\times10^{8}}{354\times10^{-9}}\
     \mathrm{J}
   = \frac{5.62\times10^{-19}}{1.602\times10^{-19}}\ \mathrm{eV}
  \tag*{(1.0)}\\
E_\nu &= 3.5\ \mathrm{eV} \tag*{(1.0)}
\end{align*}
Alternative solution:
\begin{align*}
E_\nu &= k_{\mathrm{B}}T_{\mathrm{e}}\\
  &= 1.38\times10^{-23}\times8200\ \mathrm{J}
   = \frac{1.13\times10^{-19}}{1.602\times10^{-19}}\ \mathrm{eV}
  \tag*{(2.0)}\\
E_\nu &= 0.71\ \mathrm{eV} \tag*{(1.0)}
\end{align*}
\textbf{Use of either Wien's law or $E = k_{\mathrm{B}}T$ gets full credit.
$E_\nu = 3k_{\mathrm{B}}T/2$ or similar gets no credit. Answers with
$E_\nu = 3k_{\mathrm{B}}T$ or $E_\nu = 2.7k_{\mathrm{B}}T$ also get full
credit.}
\end{description}
